\section{The reversible closure of the infinite projective space}

Let $W_{\mathrm{rev}}$ denote the smallest reversible fundamental localizer containing $B\mathbb Z/2\to *$, and let $\mathcal A_{\mathrm{rev}}$ be its pointed aspherical class. By B.4 it is accessible. The statements below use $QX=\Omega^\infty\Sigma^\infty X$ for connected pointed $X$, and $\RP^\infty_k=\RP^\infty/\RP^{k-1}$ for $k\geq1$.

\begin{claim}{C.1}{Proved}{An exact closure description and the R7 operations}
The class $\mathcal A_{\mathrm{rev}}$ is the least weak-equivalence-invariant class containing $*$ and $\RP^\infty$ and closed under pointed homotopy colimits and two out of three in homotopy cofiber sequences. It is closed under retracts, suspension and desuspension, smash product with an arbitrary pointed space, finite products, extended powers, and $\Omega^r\Sigma^r$ and $Q$ on connected spaces. Moreover
\[
 X\in\mathcal A_{\mathrm{rev}}
 \quad\Longleftrightarrow\quad
 \Omega^r\Sigma^rX\in\mathcal A_{\mathrm{rev}}
 \qquad(r\geq1).
\]
The converse with $Q$ in place of $\Omega^r\Sigma^r$ is false.
\end{claim}
\begin{proof}
Let $\mathcal B$ be the least class specified by the first sentence. A.3 constructs a reversible localizer with exactly this aspherical class, so minimality gives $\mathcal A_{\mathrm{rev}}\subseteq\mathcal B$. Conversely the cofiber criterion for $W_{\mathrm{rev}}$ supplies the two-out-of-three property, and relative Theorem A supplies homotopy-colimit closure, as checked in A.3 and B.4. Thus $\mathcal B\subseteq\mathcal A_{\mathrm{rev}}$.

We verify the additional operations rather than assuming them from R7. Retracts follow from telescopes, and suspension and desuspension from $X\to *\to\Sigma X$. For fixed $X$ in the class, the functor $X\wedge -$ preserves pointed homotopy colimits. Build an arbitrary pointed CW complex $Y$ by sphere cells, using $X\wedge S^n=\Sigma^nX$; this gives $X\wedge Y$ in the class. The cofiber sequence $X\vee Y\to X\times Y\to X\wedge Y$ then gives closure under finite products when both factors belong to the class.

If a finite group acts on a space in the class, its pointed homotopy-orbit construction is a pointed homotopy colimit and remains in the class. More generally $F_+\wedge_G X^{\wedge j}$ belongs to the class whenever $F$ is a free $G$-CW complex: filter $F$ by free cells, whose successive contributions are wedges of suspensions of $X^{\wedge j}$. This includes the extended powers and the configuration-space quotients in the little-cubes filtration. The free little-$r$-cubes construction $C_rX$ has first stage $X$ and successive quotients
\[
 C_r(j)_+\wedge_{\Sigma_j}X^{\wedge j}.
\]
Its skeletal filtration consists of cofibrations after well-pointed CW replacement. Two out of three at each stage and sequential homotopy-colimit closure give $C_rX$ in the class. May's approximation theorem, \cite[Theorem 6.1]{May1972}, states that its canonical map $C_rX\to\Omega^r\Sigma^rX$ is a weak equivalence for connected $X$, including $r=\infty$. This proves all the asserted forward closures. For the finite-$r$ converse, the adjunction triangle identity makes $\Sigma^rX$ a homotopy retract of $\Sigma^r\Omega^r\Sigma^rX$; apply retract closure and desuspend $r$ times. C.3 below gives the failure of the $Q$ converse.

The fibre-sequence clause of R7 used below also holds: if $F\to E\xrightarrow{p}B$ is a homotopy fibre sequence, $B$ is connected, and $F$ is aspherical, then $p\in W_{\mathrm{rev}}$. Express the fibration as the homotopy colimit over the fundamental infinity-groupoid of $B$ of its fibres; the fibrewise maps to points are all in the localizer, and their homotopy-colimit map is $p$. Equivalently, a category-of-simplices model and relative Theorem A prove the same assertion. Two out of three then gives $E$ aspherical if and only if $B$ is. This argument verifies the specified map, not just an abstract identification of its fibre type.
\end{proof}

\begin{claim}{C.2}{Proved}{Eilenberg--Mac Lane spaces and the sphere's infinite-loop component}
The spaces $K(\mathbb Z/2,n)$ belong to $\mathcal A_{\mathrm{rev}}$ for all $n\geq1$, and $(Q_0S^0)_{(2)}$ belongs to it. The discrete space $K(\mathbb Z/2,0)$ does not.
\end{claim}
\begin{proof}
The first space is the generator. The bar construction of a connected grouplike space is the realization of the simplicial diagram of its finite powers. C.1 gives closure under those powers and their pointed realization, so iterated bars give every $K(\mathbb Z/2,n)$, $n\geq1$. C.1 also gives $Q\RP^\infty$ in the class. Its homotopy groups are the positive homotopy groups of $\Sigma^\infty\RP^\infty$, which is $2$-local: multiplication by any odd integer is a homology equivalence of this connective spectrum and hence a stable equivalence. Thus the space is already $2$-local. The space retraction in the Kahn--Priddy theorem, in precisely the form checked in A.5, puts $(Q_0S^0)_{(2)}$ in the class. Finally $\mathcal A_{\mathrm{rev}}\subseteq\Nul$ by A.3 and minimality, and A.2 excludes the disconnected discrete two-point space.
\end{proof}

\begin{claim}{C.3}{Proved}{The stabilized finite Moore spaces become aspherical as spaces}
For every $m\geq0$, $Q(\Sigma^m\RP^2)$ belongs to $\mathcal A_{\mathrm{rev}}$, although $\RP^2$ does not.
\end{claim}
\begin{proof}
Set $E=\tau_{\geq1}\mathbb S_{(2)}$. Its zeroth infinite-loop space is $(Q_0S^0)_{(2)}$, which belongs to the class by C.2. Iterated bars therefore put
\[
 Z_r=\Omega^\infty\Sigma^rE\in\mathcal A_{\mathrm{rev}}
 \qquad(r\geq0).
\]
All these spaces are connected. The cofiber sequence of spectra $E\xrightarrow{2}E\to E/2$ is also a fibre sequence, so applying $\Omega^\infty\Sigma^r$ gives a homotopy fibre sequence $Z_r\to Z_r\to\Omega^\infty\Sigma^r(E/2)$. Its base is connected, because $E/2$ is $1$-connective. The fibre-sequence rule proved in C.1 shows that this base is aspherical.

Next smash $E\to\mathbb S_{(2)}\to H\mathbb Z_{(2)}$ with $\mathbb S/2$. The result is the fibre sequence of spectra $E/2\to\mathbb S/2\to H\mathbb Z/2$. For every $r\geq1$ its $r$th infinite-loop spaces give a homotopy fibre sequence with connected base
\[
 \Omega^\infty\Sigma^r(E/2)\longrightarrow
 \Omega^\infty\Sigma^r(\mathbb S/2)\longrightarrow
 K(\mathbb Z/2,r).
\]
The fibre and base are aspherical by the preceding paragraph and C.2. The same fibre rule gives the middle term. Since $\Sigma^\infty\RP^2\simeq\Sigma\mathbb S/2$, that term with $r=m+1$ is precisely $Q(\Sigma^m\RP^2)$. Nonmembership of $\RP^2$ follows from $\mathcal A_{\mathrm{rev}}\subseteq\Nul$ and A.4. No inference from stable membership to unstable membership has been used.
\end{proof}

\begin{claim}{C.4}{Proved}{The remaining listed asphericity tests}
Among the stunted projective spaces $\RP^\infty_k$ with $k\geq1$, exactly the case $k=1$ is in $\mathcal A_{\mathrm{rev}}$. No $\RP^{2k}$, $k\geq1$, belongs to it. For $n\geq1$, the homotopy fibre of any map $S^n\to K(\mathbb Z/2,n)$ does not belong to it; in particular this holds for the mod-$2$ fundamental-class map in the question.
\end{claim}
\begin{proof}
For $k=1$ the stunted space is the generator. For $k\geq2$, if the quotient belonged to the class, the cofiber sequence
$\RP^{k-1}\to\RP^\infty\to\RP^\infty_k$
and C.1 would put $\RP^{k-1}$ in the class. This is impossible. In even positive dimension A.4 supplies its nonzero map to the corresponding $2$-local sphere. In odd dimension $d\geq3$, the same obstruction-theoretic calculation gives
$[\RP^d,S^d_{(2)}]_*\cong H^d(\RP^d;\mathbb Z_{(2)})\cong\mathbb Z_{(2)}$.
For $d=1$, $\Map_*(S^1,S^2_{(2)})\simeq\Omega S^2_{(2)}$ is not weakly contractible. Thus none is in $\Nul$. If one uses the convention $\RP^\infty_0=\RP^\infty_+$, that case is also excluded, by disconnectedness.

Every $S^n$, $n\geq1$, is outside $\Nul$: for $n\geq2$ its localization map to $S^n_{(2)}$ is essential, and for $n=1$ use the preceding loop-space calculation. The base $K(\mathbb Z/2,n)$ is connected and is in the class by C.2. If the fibre in question were in the class, C.1 would put $S^n$ there, a contradiction. This does not require the map's fibre to be finite or simply connected.
\end{proof}

\begin{claim}{C.5}{Proved}{Local spheres and an exact local-object test}
Every $S^n_{(2)}$, $n\geq2$, is $W_{\mathrm{rev}}$-local. A space $Z$ is $W_{\mathrm{rev}}$-local if and only if it inverts the set of small arrows in B.2 for $\mathcal A=\mathcal A_{\mathrm{rev}}$ and a cardinal containing the generator. Such local spaces are closed under homotopy limits, retracts, and cotensors; they are $H\mathbb Z$-local and $B\mathbb Z/2$-null.
\end{claim}
\begin{proof}
If $f:X\to Y$ is in $W_S$, its cofiber and all suspensions of its cofiber are in $\Nul$. Use the pointed arrow $f_+:X_+\to Y_+$, whose pointed cofiber is $C(f)$. The cofiber exact sequence of mapping spaces shows that $\Map_*(\Sigma f_+,S^j_{(2)})$ is an equivalence for every $j\geq2$: its source and target are grouplike loop spaces, the fibre is null, and the next cofiber term gives surjectivity on components. Hence $\Map(f,\Omega S^j_{(2)})$ is an equivalence. Apply the unbased mapping-space functor to the EHP fibre sequence of A.1 and compare the two domains $X,Y$. Homotopy limits preserve these equivalences, proving $\Map(f,S^n_{(2)})$ an equivalence. This argument includes empty and disconnected domains without choosing a preferred component. Thus the spheres are $W_S$-local and therefore $W_{\mathrm{rev}}$-local.

The asserted set of arrows generates the localizer by B.2 and B.4. Mapping into $Z$ detects a fundamental localizer, so inverting the generators is necessary and sufficient. Mapping spaces commute with limits in their target, giving homotopy-limit and retract closure; the tensor--cotensor adjunction gives cotensors. Every reversible fundamental localizer contains $W_{\mathbb Z}$, as checked by POLES B.6 and again by the literal-axiom argument in F.2, so every one of its local objects is $H\mathbb Z$-local. The defining generator gives $B\mathbb Z/2$-nullity. Nullity alone is not substituted for the stated map-locality test.
\end{proof}

\begin{claim}{C.6}{Proved}{The finite-target locality question}
The $2$-completion of every simply connected finite CW complex is $W_{\mathrm{rev}}$-local. The same is true of its $2$-localization. The complete proof is C.8--C.9 below.
\end{claim}
The initial resolving-class approach alone did not establish this assertion. The theorem of \citet[Theorem 8 and Corollaries 9 and 11]{Strom2003} places sufficiently iterated loop spaces of finite complexes in a resolving class generated by spheres and gives an unlooped conclusion in a \emph{strong} resolving class. Local objects are not automatically closed under extensions by arbitrary fibrations. C.8 instead uses the strong resolving class defined by nullity of a fixed source, and C.9 recovers actual locality by finite smash powers and a convergent loop--suspension resolution. The distinction between nullity and locality is retained at every step.

\begin{claim}{C.7}{Refuted}{Finite fundamental group and bounded cohomology do not suffice}
A connected space with finite fundamental group and bounded integral cohomology need not be $W_{\mathrm{rev}}$-local, even when it is a finite two-dimensional complex and is $B\mathbb Z/2$-null. The finite category witnessing this assertion has $35$ objects and $159$ arrows.
\end{claim}
\begin{proof}
Let $P$ be the presentation complex
\[
 \langle a,b\mid a^5b^{-3},\ b^3(ab)^{-2}\rangle.
\]
The two cellular boundary columns in dimensions two and one are $(5,-3)$ and $(-2,1)$. Their determinant is $-1$, so $P$ is connected and integrally acyclic. Its fundamental group is the binary icosahedral group $SL_2(\mathbb F_5)$ of order $120$; the same presentation and order are recorded in \citet[Example 2.38, p.~142]{FGHHatcher2002}. We also verify the group identification independently below. In particular $P$ is not weakly contractible, while its integral cohomology vanishes in all positive degrees.

Every homology equivalence belongs to $W_{\mathrm{rev}}$. Hence $P\to *$ belongs to it. If $P$ were local, applying its mapping-space test to this arrow would make $P$ weakly contractible: more explicitly, the identity of $P$ would be homotopic to a constant map, contradicting its nontrivial fundamental group. Miller's Theorem A nevertheless makes $\Map_*(B\mathbb Z/2,P)$ weakly contractible, since $P$ is finite dimensional. Thus neither generator-nullity nor the additional fundamental-group conclusion of R11 repairs the proposed locality assertion.

Here are the complete conventions for the finite witness, so that its verification does not depend on accepting PRETOT. Write the two words as $w_0=aaaaaBBB$ and $w_1=bbbBABA$, with capitals denoting inverses. Take vertices $v,c_0,c_1$, loop edges $a,b$ at $v$, and radial edges $r_{ki}$ from $c_k$ to $v$, indexed cyclically by letters of $w_k$. For the triangle $t_{ki}$ at a positive letter $e$, set its ordered faces to $(e,r_{k,i+1},r_{ki})$; at a negative letter set them to $(e,r_{ki},r_{k,i+1})$. All these formulas satisfy the semisimplicial identities. Their realization is the specified presentation complex, by contracting one radial edge for each relation disk. Let $\mathcal P$ be its face-occurrence category: objects are simplices, and an arrow to a $d$-simplex is an increasing injection into $[d]$ with the specified source face. Composition is composition of injections. The nerve is its barycentric subdivision and has the same realization.

The program \texttt{CD\_finite\_checks.py} constructs every object, arrow and composite, checks all units and associativity triples, and computes the integral boundary matrices and Smith normal forms. It obtains $35$ objects, $159$ arrows and integral homology ranks $(1,0,0)$ with no torsion. It also computes the complete Todd--Coxeter coset table of the displayed presentation over the trivial subgroup, obtaining order $120$. Finally it enumerates all determinant-one $2$ by $2$ matrices over $\mathbb F_5$, finds matrices satisfying the two relations and generating all $120$ elements, and thereby checks that the presented group maps isomorphically onto $SL_2(\mathbb F_5)$. The exact tables, matrices and assertions are saved in \texttt{code/outputs/CD\_finite\_checks.json}.

This confirms the particular PRETOT A.4/A.5 finite construction used here. Its stronger identification with a punctured Poincar\'e homology sphere is unnecessary for the counterexample, and no such identification has been assumed.
\end{proof}

\begin{claim}{C.8}{Proved}{Looped finite targets, and an unlooped H-space case}
For every simply connected finite CW complex $K$ and every $X\in\mathcal A_{\mathrm{rev}}$,
\[
 \Map_*(X,K_{(2)})\simeq *\simeq\Map_*(X,K^{\wedge}_2).
\]
The loop space $\Omega(K^{\wedge}_2)$ is $W_{\mathrm{rev}}$-local. Moreover, if $K$ is a connected finite H-space, then $K^{\wedge}_2$ itself is $W_{\mathrm{rev}}$-local.
\end{claim}
\begin{proof}
Fix $X\in\mathcal A_{\mathrm{rev}}\subseteq\Nul$, and write
$\mathcal M_X=\{Z:\Map_*(X,Z)\simeq *\}$. This is a strong resolving class: mapping spaces preserve homotopy limits, and a fibration with null base and null fibre has null total mapping space. We apply the \emph{published} Theorem~10 of \citet{Strom2003} to the collection
\[
 \mathcal E=\{S^j_{(2)}:j\geq1\}.
\]
Suspension and smash products preserve this collection up to weak equivalence, by \cite[Propositions 6.4.2 and 6.4.6]{FGHMayPonto2012}. Its suspensions belong to $\mathcal M_X$ by the defining sphere tests. Strom's theorem therefore puts in $\mathcal M_X$ every nilpotent space having finite cone length with attaching spaces finite-type wedges of members of $\mathcal E$.

For completeness, the required finite cone-length assertion can be checked without choosing a finite model of $K$ with no $1$-cells. Collapse a maximal tree so that the connected finite CW complex $K$ has one $0$-cell. For each finite $m$, the space $V_m=\bigvee_{i=1}^m\Sigma^2K$ has a finite cell decomposition with all nonbasepoint cells of dimension at least three. Its skeletal cofiber sequences have simply connected terms and finite wedges of spheres of dimension at least two as attaching spaces. Localizing these sequences at $2$ gives cofiber sequences by \cite[Propositions 6.4.1 and 6.4.4]{FGHMayPonto2012}. It follows that $(V_m)_{(2)}$ has the finite cone length required above and hence belongs to $\mathcal M_X$. The same source's suspension and wedge comparisons give
\[
 (V_m)_{(2)}\simeq\bigvee_{i=1}^m\Sigma^2(K_{(2)}).
\]
Strom's Proposition~13, with $N=2$ and the nilpotent space $K_{(2)}$, now implies $K_{(2)}\in\mathcal M_X$. This use of a \emph{strong} resolving class concerns nullity only, and does not assert that arbitrary local objects form such a class.

Next $X$ is rationally acyclic. Indeed the pointed class of rationally acyclic spaces contains $B\ZZ/2$ and is closed under pointed homotopy colimits and two out of three for cofiber sequences, so C.1 puts all of $\mathcal A_{\mathrm{rev}}$ in it. The arithmetic square
\[
 K_{(2)}\simeq K_{\QQ}\mathop{\times}\limits^h_{(K^{\wedge}_2)_{\QQ}}K^{\wedge}_2
\]
is valid here. One can check it directly rather than assume a completion theorem for arbitrary inverse limits: the homotopy groups of $K$ are finitely generated by \cite[Theorem 4.5.2]{FGHMayPonto2012}; localization and completion identify the other groups by Theorems~6.1.2 and~11.1.2 of that source. For every finitely generated abelian group $A$, the exact sequence
\[
 0\longrightarrow A\otimes\ZZ_{(2)}\longrightarrow
 (A\otimes\QQ)\oplus(A\otimes\ZZ_2)\longrightarrow
 A\otimes\QQ_2\longrightarrow0
\]
follows from $\QQ+\ZZ_2=\QQ_2$, $\QQ\cap\ZZ_2=\ZZ_{(2)}$, and flatness of $\QQ_2$. The homotopy-pullback long exact sequence therefore identifies every homotopy group with that of $K_{(2)}$, including simple connectedness of the pullback. This proves the displayed square by Whitehead's theorem. Both rational terms are rational-homology-local, as follows either from their rational Postnikov towers or from Theorem~6.1.1 of the same source. Mapping $X$ into them gives contractible spaces. Applying $\Map_*(X,-)$ to the square thus identifies $\Map_*(X,K^{\wedge}_2)$ with $\Map_*(X,K_{(2)})$, proving the second nullity assertion.

For an unpointed map $f:A\to B$ in $W_{\mathrm{rev}}$, its based cofiber $C(f)$ lies in $\mathcal A_{\mathrm{rev}}$. The rotated cofiber sequence gives
\[
 \Map_*(\Sigma B_+,K^{\wedge}_2)\longrightarrow
 \Map_*(\Sigma A_+,K^{\wedge}_2)\longrightarrow
 \Map_*(C(f),K^{\wedge}_2)\simeq *.
\]
Hence the first arrow is an equivalence on all components. By adjunction this is precisely locality of $\Omega(K^{\wedge}_2)$.

Finally let $K$ be a connected finite H-space. Its multiplication may be chosen to have a strict two-sided unit. Multiplying a finite word in a fixed left-associated order gives a continuous map $J(K)\to K$ from the James construction: deleting any basepoint letter leaves the value unchanged. Its restriction to the one-letter stage is the identity. The James equivalence $J(K)\simeq\Omega\Sigma K$ for connected $K$ therefore makes $K$ a homotopy retract of $\Omega\Sigma K$. No associativity of its multiplication is needed for this space-level retraction. Completion preserves the two maps and their retraction identity; \cite[Proposition 11.2.4]{FGHMayPonto2012} gives
\[
 K^{\wedge}_2\quad\text{a homotopy retract of}\quad
 (\Omega\Sigma K)^{\wedge}_2\simeq\Omega((\Sigma K)^{\wedge}_2).
\]
The suspension $\Sigma K$ is a simply connected finite complex, so the last loop space is local by the preceding paragraph. Closure under retracts proves that $K^{\wedge}_2$ is local. This settles the unlooped assertion for finite H-spaces directly. The next argument removes the H-space restriction for simply connected finite targets.
\end{proof}

\begin{claim}{C.9}{Proved}{Finite simply connected targets are local after localization or completion}
For every simply connected finite CW complex $K$, both $K_{(2)}$ and $K^{\wedge}_2$ are $W_{\mathrm{rev}}$-local. In particular the simply connected finite-target question recorded in C.6 has an affirmative answer.
\end{claim}
\begin{proof}
We first prove a nullity statement for finite wedges of smash powers, then use a convergent loop--suspension resolution. This avoids interchanging completion with an infinite totalization.

Put $Z=K_{(2)}$ or $Z=K^{\wedge}_2$. For positive integers $d_1,\ldots,d_s$, let
\[
 Y=\bigvee_{j=1}^s\Sigma Z^{\wedge d_j},
 \qquad
 F=\bigvee_{j=1}^s\Sigma K^{\wedge d_j}.
\]
Thus $F$ is a simply connected finite complex. We claim that
\begin{equation}\label{eq:C-finite-smash-null}
 \Map_*(X,Y)\simeq *\quad
 \text{for every }X\in\mathcal A_{\mathrm{rev}}.
\end{equation}
If $Z=K_{(2)}$, the localization comparisons for suspension, finite wedges and smash products give $Y\simeq F_{(2)}$, so this is C.8. Suppose instead that $Z=K^{\wedge}_2$. The space $Y$ is simply connected and $2$-local, by \cite[Propositions 6.4.1, 6.4.2 and 6.4.6]{FGHMayPonto2012}. The map $F\to Y$ induced by completion is a mod-$2$ homology equivalence: completion of the nilpotent space $K$ has this property, and smash products and finite wedges preserve it, by the K\"unneth theorem over $\FF_2$. Consequently
\[
 Y^{\wedge}_2\simeq F^{\wedge}_2.
\]
Here completion inverts the indicated homology equivalence between simply connected spaces, by the characterization in \cite[Theorems 11.1.1--11.1.2]{FGHMayPonto2012}.

Apply the general arithmetic fracture theorem
\cite[Theorem 13.1.4, p.~244]{FGHMayPonto2012} to $Y$:
\[
 Y\simeq Y_{\QQ}\mathop{\times}\limits^h_{(Y^{\wedge}_2)_{\QQ}}
                  Y^{\wedge}_2.
\]
The theorem applies to all nilpotent $2$-local spaces and explicitly does not require finite generation of their homotopy groups. This generality is essential because smash products of completed spaces need not have finitely generated homotopy groups. By C.8, $X$ is rationally acyclic and $\Map_*(X,F^{\wedge}_2)\simeq *$. Both rational vertices are rational-homology-local. Mapping $X$ into the fracture square therefore proves \eqref{eq:C-finite-smash-null} also in the completed case.

The same rotated-cofiber argument used in C.8 now shows that $\Omega Y$ is $W_{\mathrm{rev}}$-local. Explicitly, if $f:A\to B$ belongs to $W_{\mathrm{rev}}$, then $C(f)\in\mathcal A_{\mathrm{rev}}$, and the fibre sequence
\[
 \Map_*(\Sigma B_+,Y)\longrightarrow
 \Map_*(\Sigma A_+,Y)\longrightarrow
 \Map_*(C(f),Y)\simeq *
\]
proves locality of $\Omega Y$ on all mapping-space components.

We next show that every positive iterate of the derived functor $T=\Omega\Sigma$ sends $Z$ to a local object. The James equivalence and splitting give, for every connected pointed CW complex $V$,
\begin{equation}\label{eq:C-James-splitting}
 \Sigma T V\simeq
 \bigvee_{q\geq1}\Sigma V^{\wedge q};
\end{equation}
see \cite[Proposition 4I.2 and Theorem 4J.1]{FGHHatcher2002}.
Inductively there are finite nonnegative integers $a_{r,m}$, for $r\geq0$ and $m\geq1$, with
\begin{equation}\label{eq:C-iterated-James}
 \Sigma T^r Z\simeq
 \bigvee_{m\geq1}\bigvee^{a_{r,m}}\Sigma Z^{\wedge m}.
\end{equation}
For $r=0$ take $a_{0,1}=1$ and $a_{0,m}=0$ for $m>1$. To pass from $r$ to $r+1$, use \eqref{eq:C-James-splitting} with $V=T^rZ$. The suspension of a smash product can be moved to each factor in succession; smash distributes over wedges. Hence a splitting of $\Sigma V$ as wedges of $\Sigma Z^{\wedge m}$ gives a splitting of $\Sigma V^{\wedge q}$ indexed by the $q$-tuples of positive weights. More explicitly,
\[
 a_{r+1,m}=
 \sum_{q=1}^m\ \sum_{\substack{m_1+\cdots+m_q=m\\m_i\geq1}}
          \prod_{i=1}^q a_{r,m_i}.
\]
This is a finite sum of finite products, which proves the asserted finiteness at every weight. The splittings need only identify the homotopy types of the individual terms; no compatibility with the later cosimplicial maps is being assumed.

Let $Y_{r,M}$ be the finite wedge on the right of \eqref{eq:C-iterated-James} with $m\leq M$, and let the maps $Y_{r,M+1}\to Y_{r,M}$ collapse the omitted summands. Since $Z$ is simply connected, $\Sigma Z^{\wedge m}$ is at least $2m$-connected. Thus the omitted tails become arbitrarily highly connected. For each fixed homotopy degree, the homotopy groups of this inverse tower are eventually constant and agree with those of the infinite wedge. The inverse-limit long exact sequence, whose $\lim^1$ term vanishes for this eventually constant tower, proves
\[
 \Sigma T^r Z\simeq\holim_M Y_{r,M}.
\]
Every $\Omega Y_{r,M}$ is local by \eqref{eq:C-finite-smash-null} and cofiber rotation. Since loops and mapping spaces preserve homotopy limits, it follows that
\[
 T^{r+1}Z\simeq\holim_M\Omega Y_{r,M}
\]
is local for every $r\geq0$.

Finally, the loop--suspension completion theorem gives, for any simply connected pointed space $Z$, the coaugmented cosimplicial resolution
\[
 Z\simeq\holim_{[q]\in\Delta} T^{q+1}Z.
\]
We use the derived resolution, so the displayed powers specify its terms up to weak equivalence. Bousfield constructs it in \cite[\S7.9]{FGHBousfield2003} and identifies its totalization with integral Bousfield--Kan completion in \S9.9; on simply connected spaces the coaugmentation is a weak equivalence. The same precise conclusion, without finiteness hypotheses, is stated explicitly in \cite[Theorem 1.8(b)]{FGHBlomquistHarper2017}. Both choices of $Z$ above are simply connected. All terms of this resolution are local by the preceding paragraph, and local objects are closed under homotopy limits. Hence $Z$ is local, proving both assertions. In particular this last step uses an ordinary loop--suspension resolution of the already completed space when $Z=K^{\wedge}_2$; it assumes no commutation of suspension or totalization with completion.
\end{proof}
